\section{\uppercase{Introduction}}
\label{sec:intro}
\noindent A longitudinal question such as ``Has my activity changed, and am I different from my peers?'' contains several analytical choices. Activity must be operationalized; earlier and recent periods must be defined; peers must be identified; and eligibility and measurement availability must be checked. The same recorded value may be close to a contemporary course mean yet below that course's earlier mean. An absent recent value may indicate administrative ineligibility, while an eligible week with no recorded clicks is a recorded zero. A dashboard that obscures these distinctions can mislead even when its displayed arithmetic is correct.\draftassist

Retrieving a useful result, selecting the intended comparison, and displaying all requested answers are different accomplishments. A source-bound validator can verify a returned result without establishing that it answers the requested window. A compiler can supply essential context that the agent did not select. Evaluating only rendering success or agreement with source numbers would conflate these responsibilities.

We investigate this distinction in educational longitudinal dashboards. A model selects predefined analyses and answers; deterministic tools calculate quantities; validation checks evidence; and a compiler constructs claims and panels. A later presentation layer adapts saved content to personal change, reference comparisons, limited observations, and assessment records. This bounded application permits studying selection failures without delegating arithmetic or unrestricted frontend code to the model.

The empirical focus is a controlled comparison of three selection interfaces using the same checkpoint and analytical backend. A compact interface is compared with derived support/scope metadata and with additional explicit analytical identity. The hypothesis is that making feature, windows, reference, and peer-versus-personal meaning explicit reduces incompleteness. Deterministic enumeration answers the same task families.

Three research questions guide the study. First, does explicit answer identity improve complete valid visible answers under fixed tool and inference budgets? Second, how do conclusions change when completeness attributable to method selections is separated from completeness supplied by compiler context? Third, what do saved failures reveal about the distinction between retrieving evidence, binding it correctly, and answering the question?

The contribution is an implemented responsibility separation, a paired development ablation, and an analysis of completeness failures and evaluation ambiguities. The negative result is informative: valid bindings do not prevent irrelevant selection, and acceptance need not imply coverage. Documented equivalent-answer forms require human adjudication. The findings do not establish agent superiority, unrestricted understanding, or a usability gain.
